\Chapter{24}

\Verse{1}
{
	\hP{}
}
{
	\eP{
		24

¹ And YHWH spoke to Moses, saying,\\
	}
}
{
}

\Verse{2}
{
	\hP{}
}
{
	\eP{
		“Command the children of Israel that they shall take clear pressed
olive oil to you for the light, to keep up a lamp always.\\
	}
}
{
}

\Verse{3}
{
	\hP{}
}
{
	\eP{
		Outside of the pavilion of the Testimony, in the Tent of Meeting, Aaron
shall arrange it from evening until morning in front of YHWH always, an
eternal law through your generations.\\
	}
}
{
}

\Verse{4}
{
	\hP{}
}
{
	\eP{
		He shall arrange the lamps on the pure menorah in front of YHWH always.\\
	}
}
{
}

\Verse{5}
{
	\hP{}
}
{
	\eP{
		“And you shall take fine flour and bake it into twelve loaves: each
loaf shall be two-tenths of a measure.\\
	}
}
{
}

\Verse{6}
{
	\hP{}
}
{
	\eP{
		And you shall set them in two rows, six to a row, on the pure table in
front of YHWH.\\
	}
}
{
}

\Verse{7}
{
	\hP{}
}
{
	\eP{
		And you shall put clear frankincense on each row, and it will be a
representative portion for the bread, an offering by fire to YHWH.\\
	}
}
{
}

\Verse{8}
{
	\hP{}
}
{
	\eP{
		On the Sabbath day, on the Sabbath day, he shall arrange it in front of
YHWH always from the children of Israel, an eternal covenant.\\
	}
}
{
}

\Verse{9}
{
	\hP{}
}
{
	\eP{
		And it shall be Aaron’s and his sons’, and they shall eat it in a holy
place, because it is the holy of holies to him from YHWH’s offerings by
fire, an eternal law.”\\
	}
}
{
}

\Verse{10}
{
	\hP{}
}
{
	\eP{
		And a son of an Israelite woman—and he was a son of an Egyptian
man—went out among the children of Israel, and the son of the Israelite
woman fought with an Israelite man in the camp.

24:10–23. a son of an Israelite woman.... they battered him with stone.
It is curious that this single story occurs here in a section fifteen
chapters long that otherwise simply states the law without any
accompanying narratives. This story’s significance may lie in the fact
that it states that the man in question is the son of an Israelite woman
and an Egyptian father, and the deity’s decision includes the notation
that the law applies equally to an alien and to a citizen. It therefore
appears that the offender in this story is an alien (see the comment on
Lev 24:16,22), and the account therefore has the value not only of
dramatizing the law concerning blasphemy, but also of dramatizing the
principle of equality before the law: “You shall have one judgment: it
will be the same for the alien and the citizen” (24:22). Indeed, the
group of laws that were described above as reflecting the notion of
equivalence of justice (“eye for eye … ”) is embedded in the middle of
this very story. The story may appear here, therefore, because it
illustrates several important laws and an extremely important legal
principle.

The story is also interesting in parallel with the only other narrative
in Leviticus, the story of the consecration of the priesthood. Like that
other story, which culminates in the deaths of Aaron’s sons, the account
of the blasphemy expresses what is at stake in the people’s young
relationship with their God. Further, in the earlier account, it is God
who deals with the offenders, Nadab and Abihu, directly; whereas in this
account of the blasphemy, humans themselves must deal justice. The two
stories thus convey together the idea that the law is both a divine and
a human concern.\\
	}
}
{
}

\Verse{11}
{
	\hP{}
}
{
	\eP{
		And the son of the Israelite woman profaned the name and cursed. And
they brought him to Moses. And his mother’s name was Shelomith, daughter
of Dibri, of the tribe of Dan.\\
	}
}
{
}

\Verse{12}
{
	\hP{}
}
{
	\eP{
		And they left him under watch, to determine it for them by YHWH’s
word.\\
	}
}
{
}

\Verse{13}
{
	\hP{}
}
{
	\eP{
		And YHWH spoke to Moses, saying,\\
	}
}
{
}

\Verse{14}
{
	\hP{}
}
{
	\eP{
		“Take out the one who cursed to the outside of the camp, and let all
who heard lay their hands on his head, and all the congregation shall
batter him.\\
	}
}
{
}

\Verse{15}
{
	\hP{}
}
{
	\eP{
		“And you shall speak to the children of Israel, saying: Any man who
curses his God: he shall bear his sin.\\
	}
}
{
}

\Verse{16}
{
	\hP{}
}
{
	\eP{
		And one who profanes YHWH’s name shall be put to death. All the
congregation shall batter him. The same for the alien and the citizen:
when he profanes the name he shall be put to death.\\
	}
}
{
}

\Verse{17}
{
	\hP{}
}
{
	\eP{
		“And a man who will strike any human’s life shall be put to death,\\
	}
}
{
}

\Verse{18}
{
	\hP{}
}
{
	\eP{
		and one who strikes an animal’s life shall pay for it: a life for a
life.\\
	}
}
{
}

\Verse{19}
{
	\hP{}
}
{
	\eP{
		And a man who will make an injury in his fellow: as he has done, so it
shall be done to him.\\
	}
}
{
}

\Verse{20}
{
	\hP{}
}
{
	\eP{
		A break for a break, an eye for an eye, a tooth for a tooth: as he
will make an injury in a human, so it shall be made in him.\\
	}
}
{
}

\Verse{21}
{
	\hP{}
}
{
	\eP{
		And one who strikes an animal shall pay for it, and one who strikes a
human shall be put to death.\\
	}
}
{
}

\Verse{22}
{
	\hP{}
}
{
	\eP{
		You shall have one judgment: it will be the same for the alien and the
citizen; because I am YHWH, your God.”\\
	}
}
{
}

\Verse{23}
{
	\hP{}
}
{
	\eP{
		And Moses spoke to the children of Israel. And they brought out the
one who cursed to the outside of the camp, and they battered him with
stone.

And the children of Israel had done as YHWH commanded Moses.\\
	}
}
{
}

